\chapter{Paged KV, prefix sharing and allocation}
\label{evod36:chapter}

Two requests consume the same three-token prompt. One continues with A and
the other with T. With two token positions per page, which bytes can they
still share? The first full page is easy: neither continuation changes it.
The second page is more subtle. It contains one valid prefix position and
one vacant slot. Writing into that vacancy through a shared pointer would
silently change the other request's history.

Chapter 35 made the ordered, per-layer Hermon cache explicit. This chapter
changes physical placement while preserving that computation. Its allocator
stores the actual two-layer Chapter 23 Transformer state through Chapter 35's
adapter, including learned absolute positions. It is a serialized CPU
reference with expensive validation and contiguous staging. It establishes
allocation and continuation contracts; it is not a production paged kernel.

\section{The logical object must survive relocation}
\index{KV cache!paged layout}\index{Logical ordinal}\index{Page table}
\label{evod36:invariant}
For layer $\ell=0,\ldots,L-1$, head count $H$ and head width $d_h$, a
consumed prefix of $T$ tokens has
\[
 K^\ell,V^\ell\in\mathbb R^{H\times T\times d_h}.
\]
One token slice contains K and V across \emph{all} layers and heads. The
reference uses $L=2,H=2,d_h=4$, eight-byte float64 elements: one such slice
occupies $2LHd_h(8)=256$ logical payload bytes. These are model tensors,
not biological molecules. The DNA vocabulary supplies token identity;
software allocation supplies the page mechanism.

Let $C_T$ denote the tuple of layer tensors plus owner, model artifact,
causal mode and ordered token metadata. The refinement obligation is
\[
 \operatorname{flatten}(\operatorname{paged}(C_T))=C_T
\]
for every layer and token ordinal, with fresh tensor storage returned by
flatten. Physical page IDs may change. Logical order, learned-position
ordinal and layer identity may not. Merely obtaining the right final
length does not establish this equality.

Choose page capacity $P$ token positions. Token ordinal $t$ uses logical
block $b=\lfloor t/P\rfloor$ and offset $o=t\bmod P$. A request's
\Term{block table} maps $b$ to a physical reference. With $P=2$, ordinal
four is block two, offset zero, even when that block resides at physical
page three. Attention must never sort page IDs to infer history.

\Plate{01-translation}{Logical ordinals pass through a block table to scattered storage. One occupied cell is the complete layer/head K/V slice at that ordinal.}{evod36:translation}

The reference deliberately groups all layers into one allocation unit. A
production engine may instead coordinate separate per-layer arenas, but
then publication must make all layer tables consistent. Allocating one
layer successfully and abandoning the others is not a valid half-cache.

\section{References carry identity; counts carry lifetime}
\index{Reference count}\index{Reference conservation}\index{ABA problem}
\label{evod36:lifetime}
A \Term{generational reference} is $\rho=(i,g)$: slot $i$ and allocation
generation $g$. Reclamation clears the page's logical payload and advances
its generation. A stale reference to $(i,g)$ fails after reuse as $(i,g+1)$.
This avoids an ABA mistake in which an old numeric slot ID appears valid
again merely because a new object occupies the slot.

The request handle additionally contains an arena-instance identity and a
monotonically issued serial. Released handles fail; a handle from another
arena fails even when slot numbers happen to match. Owner checks bind a
tenant/request pair. Refcounts alone do not supply that authorization.

Let $\mathcal S$ contain every registered live request, branch and cached
prefix table. Our central conservation invariant is
\begin{equation}
 r_i=\sum_{s\in\mathcal S}\sum_{\rho\in s.\mathrm{pages}}
       [\rho.\mathrm{id}=i],\qquad
 r_i=0\ \Longleftrightarrow\ i\text{ is free}.
 \label{evod36:counts}
\end{equation}
No sequence repeats a physical reference within its own table. A free slot
has no logical K/V tensors. A live page's generation agrees with every
table reference, and its valid row count matches the logical span.

\begin{proposition}[Reference conservation]
Assume serialized operations, valid initial registries and no process
failure during publication. Empty creation, same-tenant fresh-owner fork,
append and release preserve Equation~\eqref{evod36:counts}.
\end{proposition}
\begin{proof}
An empty table adds no references. Fork adds one table entry and increments
the corresponding count once for each entry; both sides increase equally.
Release removes those entries and decrements the same counts. If a count
becomes zero, no registered table refers to that page, so clearing it and
advancing generation cannot invalidate a live table. Append either preserves
an existing private entry, adds a new entry/count pair, or replaces one
shared tail entry by a new private entry: decrement old, increment new.
Unchanged prefix entries keep their counts. These cases exhaust the
implemented transitions. Induction over operations proves the invariant.
\end{proof}

\Listing{PagedKV.fork}

\Plate{03-lifetime}{Cache eviction drops only the cache's reference. A live child keeps the page alive until its own release. Generation changes only at zero.}{evod36:lifeplate}

The executable \texttt{audit} recomputes expected counts from the registry
rather than trusting the counters it is meant to check. A negative test
adds an unowned count and requires audit failure. A counter that is merely
nonnegative would miss that leak.
\Listing{PagedKV.audit}

\section{Append requires a partial-tail rule}
\label{evod36:cow}
If $T\bmod P=0$, all old pages are full. Append allocates a new page; it
never writes the old full page. Otherwise there is a partial tail of $o$
valid positions, where $1\le o<P$. A private tail ($r=1$) may receive new
rows. A shared tail ($r>1$) needs \Term{copy-on-write}: copy its valid
rows into private storage, redirect this branch, and append there.

For a three-token prefix and $P=2$, parent and child initially hold the
same two references. After child append, they still share the first page.
The parent tail remains one row long; the child's private tail is two rows
long. A later append by the parent may use its now-private tail without a
new allocation. Vacant capacity is not a consumed token and cannot be
included in attention normalization.

\Plate{02-cow}{Appending to a shared partial tail redirects only the child. Full pages stay shared and the parent's logical history is unchanged.}{evod36:cowplate}

\begin{proposition}[Branch isolation and native parity]
Fix the actual Hermon artifact, causal mode, learned-position convention
and deterministic evaluation. If flattened registered caches equal their
Chapter 35 counterparts, a successful append preserves that equality and
does not change any other registered sequence's logical tensors or tokens.
\end{proposition}
\begin{proof}
Chapter 35's complete-stack continuation computes the correct candidate
layer tensors from the flattened old cache. Split those candidates at
the same logical multiples of $P$. Old full blocks remain identical by
causality. New pages contain exactly the candidate slices; a private tail
has no other reader, and a shared tail is replaced only in the appending
table. Thus concatenating its page slices returns the candidate in order,
while all other tables still see their old slices. Learned positions are
derived from consumed length, independent of page ID. This proves the
step; induction over forks and continuation chunks proves the claim.
\end{proof}

The real-arithmetic statement inherits Chapter 35's assumptions. CPU
float64 tests use $2\times10^{-12}$ absolute/relative tolerance because
different chunk shapes can change rounding. This is a checked fixture
tolerance, not a numerical theorem for all Transformer artifacts.

\section{Preflight the entire chunk before publishing it}
\index{Failure atomicity}\index{Allocation!preflight}
\label{evod36:atomic}
Let $c$ be append length, $F$ the number of free pages, and $o=T\bmod P$.
For $c>0$, the required \emph{additional} page count is
\[
 N_{\rm add}=\begin{cases}
 \lceil c/P\rceil,&o=0,\\
 [r_{\rm tail}>1]+\left\lceil\max(0,c-(P-o))/P\right\rceil,&o>0.
 \end{cases}
\]
The second term subtracts the remaining capacity in the tail; the first
term pays for copying a shared tail. For $c=0$, the operation is identity
and needs zero pages. Reject if $N_{\rm add}>F$ before modifying any table,
payload, generation or refcount.

\Listing{PagedKV.required_pages}

With $T=1,P=2,c=5$ and a shared tail, one new page holds the copied old row
and first appended row; two further pages hold the other four rows.
The requirement is three. A loop that allocates one page and then discovers
exhaustion would need rollback; preflight avoids exposing that partial state.

\Plate{04-preflight}{Preflight counts COW and every subsequent block. A failed multi-page append leaves the original snapshot unchanged.}{evod36:preflight}

\Listing{PagedKV.consume}

All model arithmetic, tensor construction and content/table hashes occur
before the short publication phase. The phase installs prepared page
records and the new table without yielding. This establishes failure
atomicity for the checked validation, arithmetic and capacity paths under
the serialized model. It is not a durable transaction, a lock-free allocator
or a proof against asynchronous process interruption. Device execution
would also require pins/completion fences before reclamation, as earlier
state-pool chapters explain.

\section{Prefix caching owns pages too}
\index{Prefix cache}\index{Tenant namespace}\index{Eviction}
\label{evod36:prefix}
Deterministic execution of an identical prefix under the same artifact
produces the same per-layer K/V. The cache key therefore binds tenant,
artifact digest, causal/position mode and the exact ordered token tuple.
If adapters, attention masks or other model-affecting inputs are added,
they must enter that identity. Token text alone is insufficient.

The bounded prefix cache retains a registered fork with its own owner.
Acquiring a hit makes another fresh-owner fork; evicting releases only the
cache-owned handle. LRU orders cached entries, not physical page IDs.
The test retains a live hit, evicts its cache entry, releases its parent
and continues reading the surviving branch. Each reference remains counted.

Sharing is tenant-scoped here. An exact key and a content hash do not prove
permission to acquire another tenant's state. Cross-tenant public-prefix
sharing would need a policy and side-channel analysis: even an immutable
hit can reveal whether a prefix was present through timing. This local
Python registry is not a hostile-process security boundary.

Block-table integrity matters separately from page contents. Reordering
two valid page references changes attention history. The reference seals
ordered table metadata and layer payloads to detect accidental mutation.
It rejects a changed model artifact or stale generation. A producer able
to rewrite both data and seals can still forge a false cache; these are
consistency checks, not authentication or proof of native computation.

\section{Count physical capacity and logical sharing separately}
\index{Fragmentation!internal}\index{Physical utilization}\index{Logical multiplicity}
\label{evod36:memory}
For a nonempty unshared sequence,
\[
 N_p=\lceil T/P\rceil,\qquad W=N_pP-T,\qquad0\le W<P.
\]
Only its tail is partly unused. Fixed-size paging removes the need for a
single contiguous region at this allocation level; it does not eliminate
alignment, arena reserve or device-allocator fragmentation.

\begin{center}\input{\Generated/geometry.tex}\end{center}
\Plate{05-geometry}{The same thirteen rows use different table lengths and unused tails. Smaller pages reduce worst-case tail waste while increasing metadata and lookup operations.}{evod36:geometry}

With shared pages, sum physical occupancy once per distinct allocated page:
\[
 U_{\rm physical}=\frac{\sum_{i:r_i>0} n_i}{NP},\qquad
 M_{\rm payload\ capacity}=NPs,\qquad s=2LHd_hb.
\]
Here $n_i$ is the valid row count, $N$ is the number of allocated pages,
and $b$ is bytes per element. When $N=0$, utilization is undefined; report
zero live payload separately. The sum $\sum_r T_r/(NP)$ instead measures
logical multiplicity and can exceed one. It must not be called physical
utilization.

In our prefix-three/child-four example there are three distinct pages of
capacity two, storing five physical rows. Hence $U_{\rm physical}=5/6$.
The two logical histories contain seven rows, giving multiplicity $7/6$.
At $s=256$, physical payload capacity is 1,536 bytes, occupied rows are
1,280 bytes, and equivalent independent contiguous payload is 1,792 bytes.
Block tables, Python objects, hashes, token metadata and staging tensors
are excluded. Cached owners add references without adding physical rows.

\begin{center}\input{\Generated/branches.tex}\end{center}

For $B$ branches at a common prefix, full shared pages avoid $B-1$ copies
each. Partial-tail savings depend on which branches append and whether the
original tail remains cached. Fanout alone is not a savings percentage.
PagedAttention supplies the established serving-system comparison
for block tables and COW~\cite{evod36:kwon}; its empirical speedups are
not measurements of this reference.

\section{What the implementation establishes}
\label{evod36:evidence}
The tests compare every layer against the independent original full Hermon
path across page sizes and split positions; force scattered placement;
exercise full/partial tails, stale generations, foreign handles, cache
eviction, all branch-release orders, invalid input and multi-page OOM;
and deliberately corrupt tables, payloads and counters. These controls
make incorrect order, leaks and premature reclamation observable.

Flattening and computing through Chapter 35 still copy contiguous tensors.
Even a logically private tail is replaced by a prepared tensor object.
Thus the reference demonstrates bounded physical-page placement and sharing,
but does not claim reduced actual Python traffic, native in-place writes,
GPU coalescing, concurrent correctness or speedup. Chapter 38 supplies a
page-streaming attention reduction; Chapter 37 first treats pages as a
finite scheduling resource.

\section{Selective glossary}
\begin{description}
\item[Logical block] A consecutive span of token ordinals of capacity $P$.
\item[Physical page] An allocation unit holding all layer/head K/V at that span in this reference.
\item[Generation] Allocation identity advanced when a page is reclaimed; it is distinct from token position.
\item[Reference count] Number of registered owning table entries, including cache entries.
\item[Copy-on-write] Private replacement of a shared partial tail before an append changes its valid rows.
\item[Physical utilization] Distinct occupied page rows divided by distinct allocated page capacity.
\item[Logical multiplicity] Repeated request-visible rows per physical slot; sharing can make it exceed one.
\end{description}

\section{Twenty worked exercises}
\label{evod36:exercises}
\begin{enumerate}
\item \textbf{Translate.} Ordinal 13, $P=4$? \emph{Solution:} block 3,
offset 1; look up entry 3, never physical page 13.
\item \textbf{Size one cell.} $L=2,H=2,d_h=4,b=8$? \emph{Solution:}
$2(2)(2)(4)(8)=256$ bytes, including both K and V.
\item \textbf{Count a page.} $P=2$ for that geometry? \emph{Solution:}
512 bytes of payload capacity, before metadata.
\item \textbf{Bound waste.} Show $W<P$. \emph{Solution:}
$N_p-1<T/P\le N_p$ implies $0\le N_pP-T<P$.
\item \textbf{Empty state.} How many pages? \emph{Solution:} zero;
no all-masked pseudo-token is created.
\item \textbf{Trace COW.} Prefix three, $P=2$, fork then child append?
\emph{Solution:} first page remains shared; child gets a private second page.
\item \textbf{Full tail.} Prefix four, $P=2$, child append?
\emph{Solution:} retain both shared full pages and allocate a third.
\item \textbf{Private capacity.} One row in a private $P=2$ tail with
zero free pages can append one row? \emph{Solution:} yes, no new page credit.
\item \textbf{Preflight.} Shared $T=1,P=2,c=5$? \emph{Solution:}
$1+\lceil(5-1)/2\rceil=3$ additional pages.
\item \textbf{Evict.} Count two becomes one. Free? \emph{Solution:} no;
the remaining table still owns the page.
\item \textbf{Invalidate reuse.} Why increment generation?
\emph{Solution:} an old reference cannot name a later allocation in that slot.
\item \textbf{Audit a leak.} Count positive with no registered owner?
\emph{Solution:} recomputed incidence is zero; conservation fails.
\item \textbf{Reject a reorder.} Do page hashes alone suffice?
\emph{Solution:} no; ordered table metadata must also be checked.
\item \textbf{Prefix identity.} Same tokens under a changed position table?
\emph{Solution:} different artifact, so reject reuse.
\item \textbf{Separate permission.} Does a matching digest authorize a hit?
\emph{Solution:} no; enforce tenant/trust-domain policy first.
\item \textbf{Interpret $7/6$.} Is physical utilization above 100 percent?
\emph{Solution:} it is logical multiplicity, counting shared rows repeatedly.
\item \textbf{Explain parity.} What passes from Chapter 35 into the proof?
\emph{Solution:} exact ordered old K/V and identical learned positions
produce correct candidate rows through every layer.
\item \textbf{Inject failure.} What snapshots should an OOM test compare?
\emph{Solution:} both branches' flattened contents, tables, counts and
generations, including a failure requiring several pages.
\item \textbf{Design a page-size experiment.} \emph{Rubric:} declare
length/fanout distributions, capacity, metadata, high-water occupancy,
copy volume and end-to-end latency; do not infer speed from waste alone.
\item \textbf{Design device reclamation.} \emph{Rubric:} name in-flight
pins, completion acknowledgement, cancelled writes and stale generations;
prove reuse waits for the final physical reader, not only a logical eviction.
\end{enumerate}

\section{From storage to service}
The allocator preserves each request's causal history but cannot choose
which waiting request deserves the next page or token step. Chapter 37
adds iteration-level admission, work budgets and explicit precision
policies. These mechanisms preserve Hermon's Transformer identity;
they do not turn it into DOGMA's non-Transformer state representation.
